\begin{maintask}{4-7}
  \subsection*{課題内容}
  IPv6対応の3種類のTCPサーバ（select，fork，multithreading）を起動し，IPv6対応クライアントから接続できることを確認する．

  \subsection*{実験方法}
  \begin{enumerate}
    \item 各サーバプログラムを\texttt{~/socket}で順番に起動する．
    \item クライアントをIPv6対応にし，サーバのIPv6アドレスに接続してメッセージを送信する．
    \item サーバごとにクライアントとサーバ両方の出力を記録し，次のサーバを試す前に停止する．
  \end{enumerate}

  \subsection*{結果}
  各サーバに対するクライアント・サーバの実行結果を出力\ref{out:select_server}〜出力\ref{out:thread_client}に示す．使用したソースコードをリスト\ref{lst:select_server}〜リスト\ref{lst:ipv6_client}に示す．

  \begin{listing}[H]
    \caption{IPv6対応tcp\_server\_select.py}\label{lst:select_server}
    \begin{minted}[fontsize=\small]{python}
    #!/usr/bin/python3

    from datetime import datetime
    import socket
    import select
    import sys

    HOSTNAME = ''
    PORT_NUM = 6789
    MAX_SIZE = 1024

    print('The server started at', datetime.now())
    print('Waiting...')

    def accept_loop(server_sock):
      fds = [server_sock]
      while True:
        r, _, _ = select.select(fds, [], [])
        for sock in r:
          if sock == server_sock:
            new_sock, addr = sock.accept()
            fds.append(new_sock)
          else:
            data = sock.recv(MAX_SIZE)
            if not data:
              break
            print('Received text:', data.decode('utf-8'), 'from', addr)
            sock.sendall(data)
            sock.close()
            fds.remove(sock)

    if __name__ == '__main__':
      server = socket.socket(socket.AF_INET6, socket.SOCK_STREAM)
      #server = socket.socket(socket.AF_INET, socket.SOCK_STREAM)
      server.bind((HOSTNAME, PORT_NUM))
      server.listen()
      try:
        accept_loop(server)
      except KeyboardInterrupt:
        server.close()
    \end{minted}
  \end{listing}
  \begin{listing}[H]
    \caption{IPv6対応tcp\_server\_fork.py}\label{lst:fork_server}
    \begin{minted}[fontsize=\small]{python}
    #!/usr/bin/python3

    from datetime import datetime
    import socket
    import os

    HOSTNAME = ''
    PORT_NUM = 6789
    MAX_SIZE = 1024

    print('The server started at', datetime.now())
    print('Waiting...')

    #with socket.socket(socket.AF_INET, socket.SOCK_STREAM) as server:
    with socket.socket(socket.AF_INET6, socket.SOCK_STREAM) as server:
      server.bind((HOSTNAME, PORT_NUM))
      server.listen()
      try:
        while True:
          client, addr = server.accept()
          child_pid = os.fork()
          if child_pid == 0:
            data = client.recv(MAX_SIZE)
            if data:
              print('Received text:', data.decode('utf-8'), 'from', addr)
              client.sendall(data)
              client.close()
      except KeyboardInterrupt:
        pass
    \end{minted}
  \end{listing}
  \begin{listing}[H]
    \caption{IPv6対応tcp\_server\_multithreading.py}\label{lst:thread_server}
    \begin{minted}[fontsize=\small]{python}
    #!/usr/bin/python3

    from datetime import datetime
    import socket
    import threading

    alive = True
    observe_thread = None
    HOSTNAME = ''
    PORT_NUM = 6789
    MAX_SIZE = 1024

    print('The server started at', datetime.now())
    print('Waiting...')

    def client_handler(client, addr):
      data = client.recv(MAX_SIZE)
      if data:
        print('Received text:', data.decode('utf-8'), 'from', addr)
        client.sendall(data)
        client.close()

    if __name__ == '__main__':
      #with socket.socket(socket.AF_INET, socket.SOCK_STREAM) as server:
      with socket.socket(socket.AF_INET6, socket.SOCK_STREAM) as server:
        server.bind((HOSTNAME, PORT_NUM))
        server.listen()
        try:
          while True:
            client, addr = server.accept()
            thread = threading.Thread(target = client_handler, args = (client, addr))
            thread.start()
            while thread.is_alive():
              thread.join()
        except KeyboardInterrupt:
          alive = False
          pass
    \end{minted}
  \end{listing}
  \begin{listing}[H]
    \caption{IPv6対応tcp\_client.py}\label{lst:ipv6_client}
    \begin{minted}[fontsize=\small]{python}
    #!/usr/bin/python3

    from datetime import datetime
    import socket

    HOSTNAME = '::1'
    #HOSTNAME = '172.28.0.200'
    PORT_NUM = 6789
    MAX_SIZE= 1024

    print('The client started at', datetime.now())

    #with socket.socket(socket.AF_INET, socket.SOCK_STREAM) as server:
    with socket.socket(socket.AF_INET6, socket.SOCK_STREAM) as server:
        server.connect((HOSTNAME, PORT_NUM))
        print("Input text: ", end='')
        msg = input()
        server.sendall(msg.encode('utf-8'))
        data = server.recv(MAX_SIZE)
        print('Echo back text:', data.decode('utf-8'))
        server.close()
    \end{minted}
  \end{listing}

  \begin{termout}[H]
    \caption{selectサーバの実行結果}\label{out:select_server}
    \begin{minted}{text}
    guest@host82:~/socket$ ./tcp_server_select.py
    The server started at 2026-11-12 01:11:10.323441
    Waiting...
    Received text: How are you doing. from ('::1', 41992, 0, 0)
    guest@host82:~/socket$ ./tcp_client.py
    The client started at 2026-11-12 01:11:31.124926
    Input text: How are you doing.
    Echo back text: How are you doing.
    \end{minted}
  \end{termout}
  \begin{termout}[H]
    \caption{forkサーバの実行結果}\label{out:fork_server}
    \begin{minted}{text}
    guest@host82:~/socket$ ./tcp_server_fork.py
    The server started at 2026-11-12 01:16:08.792394
    Waiting...
    Received text: Text. from ('::1', 39216, 0, 0)
    guest@host82:~/socket$ ./tcp_client.py
    The client started at 2026-11-12 01:16:30.536484
    Input text: Text.
    Echo back text: Text.
    \end{minted}
  \end{termout}
  \begin{termout}[H]
    \caption{multithreadingサーバの実行結果}\label{out:thread_client}
    \begin{minted}{text}
    guest@host82:~/socket$ ./tcp_server_multithreading.py
    The server started at 2026-11-12 01:25:56.407150
    Waiting...
    Received text: Test message. from ('::1', 43042, 0, 0)
    guest@host82:~/socket$ ./tcp_client.py
    The client started at 2026-11-12 01:26:39.304585
    Input text: Test message.
    Echo back text: Test message.
    \end{minted}
  \end{termout}

  IPv6ループバック\texttt{::1}を用いて3種類のサーバに接続した．selectサーバでは\texttt{How are you doing.}，forkサーバでは\texttt{Text.}，multithreadingサーバでは\texttt{Test message.}を送信し，それぞれ同じ文字列がクライアントに返り，サーバ側にも接続元\texttt{::1}と受信文字列が表示された．
\end{maintask}
